\begin{abstract}
How pre-exposing an Al--Mg--Si alloy with Al$_{13}$Fe$_4$ inclusions to a
static magnetic field affects plastic deformation in creep is
examined by molecular dynamics. The inclusion is elongated along the field by
0.194\%, the strain corresponding to the 147~MPa interface stress estimated
experimentally, and is held at that strain while the matrix relaxes. The
resolved shear stress in the matrix reaches 15~MPa above the inclusion
and falls to the resolution of the calculation within 80~\AA{}.
Under shear, a pre-existing dislocation pair is torn apart from 95--105~MPa,
and no new dislocation forms at the interface up to 400~MPa in the unstrained
cell. A two-scale estimate reproduces the measured 25\% creep increase for the activation volumes reported for Al--Mg--Si,
provided the inclusions strain by 0.194\%; the magnetostriction measured for
Fe--Al is twenty times smaller, and the persistence of the effect after the
field is removed is not explained by an elastic stress.
\end{abstract}

\begin{keyword}
magnetoplasticity \sep magnetic memory \sep creep \sep molecular dynamics
\sep Al--Mg--Si alloys \sep Al$_{13}$Fe$_4$ \sep interface stress \sep
dislocations \sep solute pinning
\end{keyword}